\documentclass{article}
\begin{document}
\title{Independent chord shortening also removes the narrow first-move gap}
\maketitle

Apply the same documented waypoint-shortening control from
\#789 to the previously fixed far first approach. Conveyor or
report-side authority is unchanged. Every selected swept box
is checked against the same pinned placed collision triangles,
fences and corresponding authored gate panel. The changed
first approach is four stopped legs taking \(19.4554\) modeled
seconds, not \(26.7866\). This is a checked candidate over the
given waypoint set, not a new vehicle limit.

The courier commits far first before the signal in both worlds.
It receives the gate cue while moving on its new first diagonal
and only uses it for the return choice. On the short side
it follows the verified shortened return, and with that passage
closed it reverses the shortened first approach and uses the
other exit. The old report and old receipt, cue, new report and
new receipt use the actual pinned native broker and LUT on this
new synthetic event schedule. The cue's recipient position is
recomputed on the diagonal, not inherited from the old trajectory.
The same independently modeled single omission before enqueue
is exercised. The report remains due by 35 seconds and
no relay is registered.

With a conservative actor separation check and explicit waits,
the two modeled egresses are \(167.4460\) and \(167.0946\)
seconds or less. They meet the narrow \(167.5\)-second
diagnostic that \#789 had isolated only for the previously
frozen first-approach macros. This is not a margin for traffic or
timing uncertainty, and does not rule out roles for a timely early
cue in other physically constrained assignments.
It establishes that another ordinary control extension absorbs
this case. No independent mission, physical flight or I2 acceptance
is asserted.
\end{document}
